% !TeX program = xelatex
\documentclass[UTF8,a4paper,11pt,fontset=fandol]{ctexart}
\usepackage[top=23mm,bottom=24mm,left=26mm,right=26mm]{geometry}
\usepackage{amsmath,amssymb,amsthm}
\usepackage[unicode,hidelinks]{hyperref}
\hypersetup{pdftitle={李–刘哥德巴赫谱假设的移位素数和障碍},pdfsubject={狄利克雷 GRH 下的谱范数下界及无零区},pdfauthor={Leiming Fu}}
\ctexset{section={format=\large\bfseries\raggedright,aftername=\quad,beforeskip=2ex,afterskip=.8ex},subsection={format=\normalsize\bfseries\raggedright,aftername=\quad}}
\setlength{\emergencystretch}{1em}
\allowdisplaybreaks[1]
\newtheorem{theorem}{定理}
\newtheorem{lemma}[theorem]{引理}
\newtheorem{corollary}[theorem]{推论}
\renewcommand{\proofname}{证明}
\newcommand{\Sg}{\mathfrak S}
\newcommand{\RR}{\mathbb R}
\newcommand{\CC}{\mathbb C}
\newcommand{\Lam}{\Lambda}
\newcommand{\ind}{\mathbf 1}
\newcommand{\dd}{\,\mathrm d}
\newcommand{\GRH}{\mathrm{GRH}}
\DeclareMathOperator{\supp}{supp}
\title{\bfseries 李--刘哥德巴赫谱假设的移位素数和障碍}
\author{Leiming Fu}
\date{2026 年 11 月 26 日}
\begin{document}
\maketitle
\begin{abstract}
对李、刘定义的哥德巴赫残差算子 \(T_N\)，本文在狄利克雷广义黎曼猜想下证明，对每个固定的 \(d\ge0\)，有无穷多个偶数 \(N\) 满足
\(N\|T_N\|^2>(\log N)^d\)。因此假设 R 失败，假设 S 对最终宽度至少为 \(1/(4(\log N)^j)\) 的每个预定函数亦失败。证明从模 \(N\)、采样尺度 \(2N\) 的精确矩阵元出发。对二进区间内的素数作有限容斥，奇数阶的末层在两个移位中只出现一次，其余层的取整误差由权函数的总变差控制。普通素数定理与一个从特征和推导的平滑等差数列估计将这一偏差转移回算子。进一步，只用移位和选取模数，可预先规定一个宽度比每个固定对数负幂都小、仍有反例的函数。另一方面，存在某个可容许的 S 函数等价于 \(\|T_N\|/\Sg(N)\to0\)。本文的下界与这一极限相容。
\end{abstract}
\noindent\textbf{版本说明。} 本文保留最初的条件性表述。后续英文论文将狄利克雷 GRH 的形式化证明与反证合并，得到原命题的否定 $\neg\mathrm R$。是否存在可容许的 $h$ 使 $\mathrm S(h)$ 成立仍为开放问题。研究与写作使用了 OpenAI Codex。

\section{谱假设与结论}
以下所有对数均为自然对数。令 \(\Lam\) 为冯·芒戈尔特函数，\(\varphi\) 为欧拉函数，并设
\[
 C_2=\prod_{p>2}\left(1-\frac1{(p-1)^2}\right),\qquad
 \Sg(N)=2C_2\prod_{\substack{p\mid N\\p>2}}\frac{p-1}{p-2},
 \qquad \delta_N^2=\frac{\Sg(N)(\log N)^3}{N}.
\]
乘积中的 \(p\) 遍历素数。称 \(N\) 满足 \(\GRH(N)\)，若导子整除 \(N\) 的每个本原狄利克雷特征的 \(L\) 函数在开临界带中的零点实部均为 \(1/2\)。本文以 \(G\) 表示对所有正模数的狄利克雷 GRH。这与李、刘定理 1.1 和 1.2 所用的“每个充分大的偶数满足 \(\GRH(N)\)”等价。事实上，每个正导子 \(d\) 都整除任意大的偶数 \(2kd\)，故后一个条件已包含该导子的全部本原特征。李、刘在 Hardy 条带 \(\alpha_N=1/\log N<\Re s<\beta=1/8\) 上定义紧自伴残差算子 \(T_N\)，主定理使用原文归一化 \(c=1\) \cite{LL}。其非零特征值 \(\mu_\nu\) 按重数计，四元数谱函数为
\[
 X_N(q)=\prod_\nu\left(1-\frac{4\mu_\nu^2}{\delta_N^2}\,q(1-q)\right).
\]
假设 R 要求每个充分大且满足 \(\GRH(N)\) 的偶数 \(N\) 都有 \(X_N(q)=0\Rightarrow\Re q=1/2\)。对于预先给定的正函数 \(h(N)\to\infty\)，置
\[
 \eta_h(N)=\frac{h(N)^2(\log N)^3}{N\Sg(N)} .
\]
若最终 \(\eta_h(N)\le1/4\)，称 \(h\) 可容许。假设 \({\rm S}(h)\) 要求在上述偶模数上，\(\Re q>1-\eta_h(N)\) 蕴含 \(X_N(q)\ne0\)。全文的最终性和极限均沿偶数理解。

\begin{theorem}\label{th:main}
设 \(G\) 成立。对每个整数 \(d\ge0\) 和每个 \(K\)，存在偶数 \(N\ge K\) 使
\[
                    N\|T_N\|^2>(\log N)^d.                  \tag{1}
\]
故 R 不成立。若对某个固定整数 \(j\ge0\)，可容许的 \(h\) 最终满足
\(\eta_h(N)\ge 1/(4(\log N)^j)\)，则 \({\rm S}(h)\) 亦不成立。更具体地，对每个 \(j\ge0\) 和任意 \(K\)，原谱函数在某个偶数 \(N\ge K\) 处有实零点 \(x\) 满足
\[
            1-\frac{1}{4(\log N)^j}<x<1.
\]
\end{theorem}

例如，对 \(N>1\) 预先指定
\[
              h_j(N)=\sqrt{\frac{N\Sg(N)}{4(\log N)^{j+3}}}.
                                                               \tag{2}
\]
在 \(C_2\) 的乘积中写 \(n=p-1\)。补入所缺的整数 \(n\ge2\) 只会减小乘积，而
\(\prod_{n=2}^{m}(1-n^{-2})=(m+1)/(2m)\)。故 \(C_2\ge1/2\)，\(\Sg(N)\ge1\)。于是每个 \(h_j\) 都趋于无穷，其宽度恰为 \(1/(4(\log N)^j)\)。这给出各自有无穷多个反例的显式函数族。Yilin Liu 的预印本采用相反方向的素数移位与 Maier 矩阵，以 Buchstab 密度得到同类谱障碍 \cite{YL}。本文沿用检测这两个移位的思路，重新推导矩阵元，并以有限层容斥直接构造所需的算术偏差。

\section{从残差算子提取一对移位}
若 \(1\le H\le P<N/20\)，称实值 \(b\in C^\infty_c(\RR)\) 为测试权，若其支撑包含在 \([P,2P]\) 内且包含在某个长为 \(H\) 的区间内，并满足
\(\|b\|_\infty\le1,\ H\|b'\|_\infty\le1\)。定义
\[
 B_N(b)=\sum_{\substack{m<N\\(m,N)=1}}\Lam(m)b(m)
                      \{\Lam(N-m)-\Lam(N+m)\}.
\]
为明确算子，记 \(L_\chi(s)= (L'/L)(1/2+s,\chi)\) 及 \(w=s+z\)。解析残差为
\[
 \mathcal E_N(s,z)=\frac{2}{\varphi(N)}
 \sum_{\chi(-1)=-1} L_\chi(s)\overline{L_\chi(\bar z)}
 -\sum_{(n,N)=1}\frac{\Lam(n)^2}{n^{1+w}}
 -N^{-w}\mathcal D_N(w)\mathcal B(s,z).
\]
这里 \(n\) 取正整数，
\(\mathcal D_N(w)=\Sg(N)\zeta(w)\prod_{p>2,\,p\nmid N}(1+p^{-w}/(p-2))\)；
在绝对收敛域它等于 \(\sum_{\ell\ge1}\Sg(\ell N)\ell^{-w}\)。用 Euler 的 Beta 函数记
\[
 \mathcal B(s,z)=B(1/2-s,w)+B(1/2-z,w)-B(1/2-s,1/2-z).
\]
在 Hardy 条带上，\(T_N\) 的核为 \(\mathcal E_N(s,\bar z)/((1+s)(1+\bar z))\)。其作用取左右两条边界的迹与此核积分之和，边界高度取 Lebesgue 测度。下文的向量范数与积分均采用这一归一化。

\begin{lemma}\label{le:transfer}
存在绝对常数 \(C>0\)，使得对每个满足 \(\GRH(N)\) 的偶数 \(N>e^{10}\) 和每个测试权均有
\[
                     H\,B_N(b)^2\le C P^2 N\|T_N\|^2.         \tag{3}
\]
\end{lemma}
\begin{proof}
取 \(M=2N\)，令 \(f(u)=b(Mu)\)。取固定光滑函数 \(g\)，关于 \(1/2\) 对称，在 \([7/16,9/16]\) 上为一，支撑包含在 \([3/8,5/8]\)，且 \(0\le g\le1\)。此时 \(\{u:f(u)\ne0\}\subset (0,1/8)\)。对支撑在 \((0,1)\) 内的实光滑权 \(v\)，置
\[
 w_v(n)=M^{-1/2}\Lam(n)v(n/M),\qquad
 J_v(y)=2\int_{\RR}v(u)v(u+y)\dd u-\int_{\RR}v(u)v(y-u)\dd u .
\]
函数按零延拓。令 \(\widetilde v(u)=\sqrt u\,v(u)\)（\(u>0\)）。用 \(\widetilde v\) 的 Mellin 变换在尺度 \(M\) 上构成 Hardy 线性泛函，记其 Riesz 向量为 \(x_v\)。将李--刘残差核的奇特征项、精确对角项和 Beta 项分别配对，可得
\begin{align}
 \langle T_N x_v,x_v\rangle
 ={}&\frac{2}{\varphi(N)}\sum_{\chi(-1)=-1}
     \left|\sum_{n<M}w_v(n)\chi(n)\right|^2
       -\sum_{\substack{n<M\\(n,N)=1}}w_v(n)^2 \notag\\[-2pt]
    &-\sum_{\ell\ge1}\Sg(\ell N)J_v(\ell/2).                  \tag{3a}\label{eq:quad}
\end{align}
这里特征的模数仍为 \(N\)。为说明半整数采样，在两条 \(\Re s=\Re z=1/9\) 的积分线上令 \(w=s+z\)。Beta 核在内层的 Mellin 积分等于 \(\mathcal M J_v(w)\)。外层比尺度 \(N\) 的公式多出因子 \(2^w\)。在密度狄利克雷级数绝对收敛的右侧直线上，逐项反演给出 \(J_v(\ell/2)\)。移线使用正则化乘积，密度在 \(w=1\) 的极点由
\(\mathcal M J_v(1)=\int_0^\infty J_v(y)\dd y=0\) 消去；此等式由两次积分及交换 \(u\) 与 \(u+y\) 得到。光滑紧支撑给出水平边界衰减。奇特征项则由 \(-L'/L(1/2+s,\chi)\) 的 Mellin 反演得到；因 \(\GRH(N)\)，非主特征移线时无极点。对角项直接反演。式 \eqref{eq:quad} 为双线积分的精确等式。

对 \(v=f+g,f,g\) 作极化。对角交叉项由支撑分离消失。对 \(y=\ell/2\)，恒有 \(g(u)f(u+y)=0\)；当 \(f(u)\ne0\) 时，\(g(u+y)=g(y-u)\)，其中 \(\ell=1\) 用对称性，\(\ell\ge2\) 时两端均为零。反射卷积的另一个交叉项经 \(u\mapsto y-u\) 与第一个相等。因此整个 Beta 交叉项为零。奇特征正交关系给出同余类 \(n\equiv m\pmod N\) 与 \(n\equiv-m\pmod N\) 的指示函数之差，并同时要求样点为单位。在 \(g\) 的支撑中，贡献只来自 \(n=N+m,N-m\)；两处 \(g=1\)。故极化所得的实矩阵元为 \(-B_N(b)/M\)。

泛函范数的精确积分给出
\[
 \|x_v\|^2\le\frac{e^4}{2\pi}
       \int_0^\infty\left(v(u)/2-u v'(u)\right)^2\dd u .
\]
事实上，精确被积式还含有因子
\((Mu)^{2\alpha_N}/(1+(Mu)^{-2(\beta-\alpha_N)})\)，它在有效支撑上至多为 \(e^4\)。对 \(f\) 换元 \(y=Mu\) 得到因子 \(M^{-1}\)。由测试权的两种支撑约束，\(H\int (b(y)/2-yb'(y))^2\dd y\le 9P^2\)。对固定的 \(g\)，积分为常数。代入
\(|B_N(b)|/M\le\|T_N\|\|x_f\|\|x_g\|\) 即得 (3)。
\end{proof}

\section{奇数阶容斥所产生的偏差}
固定奇整数 \(k\ge1\)。令整数 \(t\) 趋于无穷，并设
\[
 \begin{gathered}
 z=8t^4,\qquad H=t^{4k},\qquad D_k=32\cdot16^k,\qquad P=D_kH,\\
 \mathcal P=\{p:z<p\le2z,\ p\ \text{为素数}\},\qquad
 s=|\mathcal P|,\qquad Q=2\prod_{p\in\mathcal P}p .
 \end{gathered}
\]
取两个固定的非负光滑函数 \(u_0,v_0\)，值域在 \([0,1]\) 中。令 \(u_0\) 在 \([2,3]\) 上为一，在 \((1,4)\) 外为零；令 \(v_0\) 在 \([7/5,8/5]\) 上为一，在 \((1,2)\) 外为零。写
\[
 b_a(x)=u_0((x-a)/H)v_0(x/P),\qquad
 R(a)=\sum_{P\le m\le 2P}\Lam(m)b_a(m)
          \{\ind_{(a-m,Q)=1}-\ind_{(a+m,Q)=1}\}.
\]
负数的最大公因数按绝对值定义。存在与 \(a,H,P,k\) 无关的 \(V\ge1\)，使 \(b_a/V\) 在 \(P<N/20\) 时是宽度 \(3H\) 的测试权。这由乘积求导以及 \(3H\le P\) 得出。

\begin{lemma}\label{le:sieve}
存在只依赖 \(k\) 的常数 \(c_k>0\)。每个充分大的整数 \(t\) 都对应某个 \(0\le a\le2P\) 的偶整数，使得 \(R(a)\ge c_k t^{3k}\)。
\end{lemma}
\begin{proof}
记 \(L=4H+1\)。对固定的奇数 \(m\in[P,2P]\)，用 \(h=m-a\) 编列，奇数 \(0<h<L\) 对应偶数 \(a\)。置 \(w(h)=u_0(h/H)\)，置
\(v(m)=\ind_{m\ {\rm odd}}\Lam(m)v_0(m/P)\)。由于 \(2\mid Q\)，对偶数 \(a,m\) 两个单位指示函数均为零。换序后，所有偶数行的和为
\[
 \sum_{P\le m\le2P}v(m)\sum_{\substack{0<h<L\\h\ {\rm odd}}}w(h)
       \big\{\ind_{(h,Q)=1}-\ind_{(2m-h,Q)=1}\big\}.          \tag{4}
\]
在权支撑上 \(h\le m\)，故换序恰好覆盖 \(0\le a\le2P\)。

对 \(\mathcal P\) 完全展开容斥。对每一子集 \(T\) 记 \(d_T=\prod_{p\in T}p\)。充分大的 \(t\) 满足 \(L\le z^k\) 和 \(4P<z^{k+1}\)。于是 \(k\) 个带内素数之积不能整除 \(h\)，且 \(k+1\) 个之积不能整除 \(2m-h\)。当 \(|T|=k\) 时，只保留 \(d_T\mid 2m-h\)；由 \(k\) 为奇数，这一项在差中系数为正一。

当 \(|T|<k\) 时，两侧分别在模 \(2d_T\) 的剩余类 \(d_T\) 和 \(2m-d_T\) 中计数。在任意前缀中，两类计数之差的绝对值至多为二。有限 Abel 求和给出每个子集、每列 \(m\) 的加权误差上界
\[
   2\left(|w(L)|+\sum_{0\le h<L}|w(h)-w(h+1)|\right)\ll 1. \tag{5}
\]
这里离散总变差由 \(\|u_0'\|_\infty L/H\) 控制。\(|T|<k\) 的子集不超过 \(k(s+1)^{k-1}\) 个；又有 \(\sum_m v(m)\le(P+1)\log(2P)\)。这同时估计了 (4) 中所有低阶项。

对每个最高层子集，\(d=d_T\le(2z)^k=P/32\) 为奇数。在到 \(3P\) 的距离不超过 \(d\) 处取 \(d\) 的奇倍数 \(y\)。若素数 \(m\) 属于
\[
              (y+2H)/2<m\le(y+3H)/2 ,
\]
则 \(h=2m-y\) 为奇，\(w(h)=1\)，且 \(7P/5\le m\le8P/5\)，因此 \(v(m)=\log m\)。区间长为 \(H/2=P/(2D_k)\)。普通素数定理 \(\vartheta(x)\sim x\) 给出其中素数对数权之和至少为 \(H/4\)，并对所有 \(d\) 一致。确实，取任意充分小的固定 \(\varepsilon_k>0\)，在 \([P,2P]\) 内用 \(|\vartheta(x)-x|\le2\varepsilon_k P\) 估计两个端点即可。把最高层逐个相加得到
\[
 \sum_{\substack{0\le a\le2P\\ a\ {\rm even}}}R(a)
 \ \ge\ \binom{s}{k}\frac H4-C k(s+1)^{k-1}(P+1)\log(2P).  \tag{6}
\]
普通素数定理在 \((z,2z]\) 上还给出 \(s\log(2z)\ge z/2\)，因而最终 \(s\ge t^3\)。对固定的 \(k\)，有 \(\log(2P)\le t\) 最终成立。当 \(s\ge2k\) 时，\(\binom{s}{k}\ge(s/2)^k/k!\)。故 (6) 的误差至多为主项的一半，只需令 \(s\) 超过依赖于 \(k\) 的常数乘 \(\log(2P)\)。除以至多 \(2P+1\) 个行指标，所得某行至少为常数倍的 \(s^kH/P\gg_k t^{3k}\)。
\end{proof}

\section{从剩余类平均回到素数移位}
令 \(g\) 为引理 \ref{le:transfer} 的平滑宽权。对正整数 \(M\) 及模 \(Q\) 的剩余类 \(v\)，置
\[
 A_{Q,M}(v)=\sum_{\substack{n<M\\n\equiv v\ (Q)}}\Lam(n)g(n/M),
 \qquad U=\frac{1}{\varphi(Q)}
       \sum_{\substack{n<M\\(n,Q)=1}}\Lam(n)g(n/M).
\]
\begin{lemma}\label{le:ap}
若 \(Q\ge4\) 满足 \(\GRH(Q)\)，则对所有可逆类 \(v\) 有
\[
                  |A_{Q,M}(v)-U|\ll_g M^{3/4}Q^2.            \tag{7}
\]
若 \(4Q\le M\) 且 \(M\) 充分大，则 \(U\ge M/(16\varphi(Q))\)。
\end{lemma}
\begin{proof}
对非主特征 \(\chi\bmod Q\)，先将
\(\sum_n\Lam(n)\chi(n)n^{-1/2}F(n/M)\) 写为
\(M^s\mathcal MF(s)(-L'/L)(1/2+s,\chi)\) 在 \(\Re s=1\) 上的 Mellin 积分。取 \(F(u)=\sqrt{u}\,g(u)\)。局部 GRH 给出被积函数在 \(1/4\le\Re s\le1\) 上的全纯性。该闭条带上，特征的对数导数有标准的 \(\ll\log^2(Q(2+|\Im s|))\) 型界；这一步使用原算子形式化中已证明的对数导数估计。四次分部积分给出 \(\mathcal MF(s)\ll_F(1+|\Im s|)^{-4}\)。把对数界粗化为 \(O(Q^2(1+|\Im s|)^2)\)，可将积分移至 \(\Re s=1/4\)，得到 \(O_g(M^{1/4}Q^2)\)。乘以 \(\sqrt M\) 即为未归一化特征和的估计。特征正交关系将非主特征的和除以 \(\varphi(Q)\)，从而得到 (7)。区间 \((7M/16,9M/16]\) 中的每个素数皆超过 \(Q\)，且 \(g=1\)。该区间的 \(\vartheta\) 增量至少为 \(M/16\)，这给出第二项结论。
\end{proof}

对 \(m\le2P\) 分别取符号 \(\pm\)，以 \(N\equiv a\bmod Q\) 为中心求和，另保留 \((N,m)=1\)。当 \(8P\le M\)、\(4Q\le M\) 且 \(m>0\) 时有
\begin{align}
 &\sum_{\substack{N<M,\;N\equiv a\ (Q)\\(N,m)=1}}
            g(N/M)\Lam(N\pm m)
       = U\ind_{(a\pm m,Q)=1}+O(E),                          \tag{7a}\label{eq:apshift}\\
 &E=C M^{3/4}Q^2+2\sqrt M\log M+D(2P)\log(2M). \notag
\end{align}
这里 \(D\) 是只依赖于 \(g\) 的常数。负移位在可能贡献的区间内恒为正。为验证误差，先将 \(g(N/M)\) 移至 \(g((N\pm m)/M)\)。按均值定理，每项权差至多为 \(\|g'\|_\infty m/M\)；以 \(\Lam(n)\le\log(2M)\) 对整个区间估计即可。支撑远离端点，换指标无边界项。设 \(n=N\pm m\)，则 \((n,m)=(N,m)\)。若 \(n\) 是贡献区中的素数，则 \(n>m,Q\)，满足两项互素限制。在可逆类，去掉的项仅来自真素数幂；在不可逆类，整个和仅来自真素数幂。两种情形各自总误差至多为 \(\psi(M)-\vartheta(M)\le2\sqrt M\log M\)。将其与 (7) 合并即得 \eqref{eq:apshift}。

取 \(r=2^{t^9}\)、\(M=r^{64}\)。对 \(t\) 充分大有 \(Q,P\le r\)。例如 \(|\mathcal P|\le 2z+1\) 且 \(p\le 2^{2z}\)，故 \(\log_2 Q\le1+2z(2z+1)\le t^9\)。固定多项式 \(P\) 亦满足 \(P\le r\)。两种符号的误差总计至多为
\[
            2E(P+1)\log(2P)\le M/(32Q).                      \tag{8}
\]
可直接逐幂验证。最终 \(C,D,\log(2M)\le r\)，因而 \(E\le5r^{51}\)；把左端乘以 \(32Q\) 至多得 \(640r^{54}\le r^{64}\)。

对引理 \ref{le:sieve} 选中的偶行 \(a\)，在 \eqref{eq:apshift} 中两式相减，乘以 \(\Lam(m)b_a(m)\) 再对 \(m\) 求和。利用 \(g\) 的支撑，在有贡献的中心有 \(2P<N\)。由引理 \ref{le:ap} 和 (8) 得
\[
 \sum_{\substack{N<M\\N\equiv a\ (Q)}}g(N/M)B_N(b_a)
     \ge U R(a)-\frac{M}{32Q}
     \ge\frac{M}{32Q}R(a) .
\]
最后一步取 \(t\) 足够大使 \(R(a)\ge1\)。可用的中心数至多为 \(M/Q+1\le2M/Q\)。因此存在偶数 \(M/4<N<M\)，在此处
\(B_N(b_a)>R(a)/128\)。对充分大的 \(t\)，有 \(80P<M\)，故 \(P<N/20\)。对归一化权 \(b_a/V\) 应用 (3)，注意其宽度为 \(3H\)，得到
\[
               N\|T_N\|^2\gg_k \frac{H t^{6k}}{P^2}
                     \gg_k t^{2k},\qquad
               \log N\le 64\,t^9\log 2 .                       \tag{9}
\]
这些中心随 \(t\) 无界。任给固定 \(d\)，先取奇数 \(k\) 使 \(2k>9d\)，再令 \(t\) 趋于无穷，即得 (1)。

\section{从范数下界到谱零点}
李--刘的谱字典可由每个特征值对应的二次因子读出。紧自伴性保证当 \(\|T_N\|>0\) 时有非零特征值达到该绝对值；平方可和性保证收敛乘积在各因子均非零时不消失。在任一复切片，二次因子的零点是
\[
          \frac{1}{2}\left(1\pm\sqrt{1-\delta_N^2/\mu_\nu^2}\right).
\]
因此所有零点实部为 \(1/2\) 等价于 \(\|T_N\|\le\delta_N\)。若 \(0<\eta<1/2\)，则半空间 \(\Re q>1-\eta\) 无零点等价于
\[
                4\eta(1-\eta)\|T_N\|^2\le\delta_N^2 .          \tag{10}
\]
边界取等与半空间严格不等号相配。若违反 (10)，上式取正号的根是落在 \((1-\eta,1)\) 内的实零点。

\begin{proof}[定理 \ref{th:main} 的证明（续）]
设 \(\omega(N)\) 为不同素因子的个数。将奇素因子依次写作 \(p_1<\cdots<p_\ell\)。因 \(p_i\ge i+2\)，有
\(\prod_i(p_i-1)/(p_i-2)\le\prod_i(i+1)/i=\ell+1\)。
又 \(2^{\omega(N)}\le N\)，故最终
\(\Sg(N)\le 2(1+\log N/\log 2)\le(\log N)^2\)。
R 会给出 \(N\|T_N\|^2\le\Sg(N)(\log N)^3\le(\log N)^5\)，与 (1) 矛盾。
若预定条带最终至少宽 \(\eta_0=1/(4(\log N)^j)\)，则 \({\rm S}(h)\) 蕴含 (10) 在 \(\eta=\eta_0\) 处成立。利用 \(1-\eta_0\ge 3/4\) 得
\(N\|T_N\|^2\le (4/3)(\log N)^{j+5}\)，与 (1) 在 \(d=j+6\) 的结论矛盾。
\end{proof}

\begin{corollary}\label{co:seq}
设 \(G\) 成立。存在同一列严格递增的偶数 \(N_i\) 和实零点 \(X_{N_i}(x_i)=0\)、\(x_i<1\)，使每个固定整数 \(j\ge0\) 都满足
\[
             (\log N_i)^j(1-x_i)\longrightarrow 0.
\]
若一个预定函数 \(h\) 满足 \({\rm S}(h)\)，则沿这同一列最终有
\(\eta_h(N_i)\le 1-x_i\)。
\end{corollary}
\begin{proof}
第 \(i\) 步应用定理 \ref{th:main} 的实零点结论，取指数 \(i\) 且要求 \(N_i>N_{i-1}\)。固定 \(j\) 后，当 \(i\ge j+1\) 且 \(\log N_i\ge1\) 时，乘积 \((\log N_i)^j(1-x_i)\) 介于零和 \(1/(4\log N_i)\) 之间。最后一个不等式由原无零区的开边界和 \(X_{N_i}(x_i)=0\) 直接得到。
\end{proof}

\section{一个比任意固定对数负幂更窄的反例}

可以在定义 \(h\) 时完全避开谱数据。对整数 \(d\ge0\)，称偶数 \(N\) 满足条件 \(A_d\)，若存在第 2 节的测试权 \(b\) 和参数 \(H,P\) 使
\[
                H B_N(b)^2>P^2(\log N)^d .
\]
设 \(r_{-1}=2\)。对 \(j\ge0\)，令 \(r_j\) 为大于 \(r_{j-1}\) 的最小偶数 \(N\)，满足 \(A_{j+8}\) 且对一切整数 \(n\ge N\) 均有 \(h_j(n)\ge j\)，其中 \(h_j\) 见 (2)。置
\[
 a(N)=\max\big(\{0\}\cup\{j:r_j\le N\}\big),\qquad
                    h_*(N)=h_{a(N)}(N).
\]
函数的定义只涉及冯·芒戈尔特函数、奇异级数和实测试权。

\begin{theorem}\label{th:arithwidth}
设 \(G\) 成立。上述每个最小值均存在；\(h_*\) 可容许，且每个固定 \(d\ge0\) 都有
\[
             (\log N)^d\eta_{h_*}(N)\longrightarrow0 .
\]
尽管如此，\({\rm S}(h_*)\) 失败。
\end{theorem}
\begin{proof}
在 (9) 的构造中，应用算子估计以前已得到 \(HB_N(b)^2/P^2\gg_k t^{2k}\)，常数吸收测试区间宽度的固定因子 \(3\)。由 \(\log N\le64t^9\) 及 \(k\) 可任意大，\(A_d\) 对每个固定 \(d\) 都在任意大的偶数上成立。又 \(h_j(n)\to\infty\)，故可依次取最小值 \(r_j\)。由严格递增性，\(a(r_j)=j\)，且 \(a(N)\to\infty\)。当 \(N\ge r_0\) 时，\(h_*(N)\ge a(N)\)。由 (2)，其宽度是 \(1/(4(\log N)^{a(N)})\)，最终至多 \(1/4\)。对固定 \(d\) 取 \(a(N)\ge d+1\) 即得所述极限。

若 \({\rm S}(h_*)\) 成立，在充分大的 \(N=r_j\) 上，(10) 及 \(\Sg(N)\le(\log N)^2\) 给出 \(N\|T_N\|^2\le(\log N)^{j+6}\)。这里的阈值与 \(j\) 无关，因为 \(1-\eta_{h_*}\ge3/4\)。将 (3) 用于 \(A_{j+8}\) 的测试权，得到 \((\log N)^{j+8}< C(\log N)^{j+6}\)，与 \(r_j\to\infty\) 矛盾。
\end{proof}

\section{慢增长函数的判别}
\begin{theorem}\label{th:slow}
在充分大的、满足 \(\GRH(N)\) 的偶数上，存在一个可容许函数 \(h\) 使 \({\rm S}(h)\) 成立，当且仅当
\[
                 \frac{\|T_N\|}{\Sg(N)}\longrightarrow 0.    \tag{11}
\]
这里的等价命题本身不需要 \(G\)。具体而言，任给 \(\varepsilon>0\)，存在 \(K_\varepsilon\)，使每个满足条件的 \(N\ge K_\varepsilon\) 均有 \(\|T_N\|\le\varepsilon\Sg(N)\)。等价地，所有可容许的 \(h\) 都使 \({\rm S}(h)\) 失败，当且仅当存在固定的 \(\varepsilon>0\) 及任意大的满足局部 GRH 的偶数 \(N\)，使 \(\|T_N\|>\varepsilon\Sg(N)\)。
\end{theorem}
\begin{proof}
代入 \(\eta_h\)，式 (10) 化为
\[
         2h(N)\sqrt{1-\eta_h(N)}\,\|T_N\|\le\Sg(N).          \tag{12}
\]
若 \({\rm S}(h)\) 成立，根号最终至少为 \(\sqrt{3}/2\)，而 \(h(N)\to\infty\)，故有 (11)。反向把 \(a_N=\|T_N\|\) 在不满足局部 GRH 或条带参数无效的指标处延拓为零。对 \(N\ge4\) 置
\[
 h(N)=
 \begin{cases}
 \min\{\log N,\Sg(N)/(2a_N)\},& a_N>0,\\
 \log N,&a_N=0.
 \end{cases}
\]
(11) 蕴含对每个固定 \(B>0\)，两个参与取最小值的量最终都至少为 \(B\)。故 \(h\) 正且发散。由 \(h(N)\le\log N\) 及 \(\Sg(N)\ge1\)，有 \(\eta_h(N)\le(\log N)^5/N\le1/4\) 最终成立。构造给出 \(2h(N)a_N\le\Sg(N)\)，从而满足 (12)。
\end{proof}

定理 \ref{th:slow} 中的一次选取是对整列 \(N\) 选择一个函数。对于指定函数 \(h(N)=\log N\)，(12) 所需的范数规模是 \(\Sg(N)/(2\log N\sqrt{1-\eta_h(N)})\)。为检验 (1) 对此有多大约束，可取标量列 \(m_N=1/(\log N)^2\)。任给 \(d\)，最终 \(Nm_N^2>(\log N)^d\)，同时 \(2(\log N)m_N\le\Sg(N)\) 最终成立。此列只用来比较两个估计的强度。

原算子还有一个可检验的正向估计。由奇特征项、精确对角项及 Beta 模型各自的 Hilbert--Schmidt 范数界相加得 \(\|T_N\|\le C(\log N)^2\)，其中只用局部 \(\GRH(N)\) \cite[定理~6.2]{LL}。在 (10) 中取 \(\eta=c\,\Sg(N)/(N\log N)\)，只要 \(0<c\le1/(4C^2)\)，即可保证该半空间内无零点。这给推论 \ref{co:seq} 中的实际根提供了另一侧的估计 \(1-x_i\ge c\,\Sg(N_i)/(N_i\log N_i)\)。对 \(h=\log\)，所要求的宽度与上述保证宽度之比为
\[
       \frac{(\log N)^6}{c\,\Sg(N)^2}\ \ge\ \frac{(\log N)^2}{c}
       \ \longrightarrow\ \infty .
\]
因此现有上界还缺少抵消三个算子分量所需的定量估计。原算子能否满足 (11)，或能否满足 \({\rm S}(\log)\)，仍未判定。李--刘的定理 1.1 表明，在 \(G\) 下证明其中任意一个相应的 S 断言，便可得到原文的加权二元哥德巴赫渐近式 \cite{LL}。

\clearpage
\section*{形式化核验}
上述有限筛、特征和平均、Mellin 配对以及谱结论在 Lean 4.35.0 中连接至李--刘项目所定义的 \texttt{residualOperator}、\texttt{HypothesisR} 和 \texttt{HypothesisS} \cite{LLLean}。定理 \ref{th:main} 的形式命题显式量化 \(G\)；另一个入口直接使用原文的“充分大偶数满足局部 GRH”条件。定理 \ref{th:slow} 的形式命题保留每个模数自己的 GRH 条件。普通素数定理使用已证明的 \(\vartheta(x)/x\to1\)。移位和偏差由有限容斥给出，范数下界由精确矩阵元推得。覆盖上述结果的同一份导出由 Lean 内核、lean4lean、con-leche 与 nanoda 分别重放。导出仅含 \texttt{propext}、\texttt{Classical.choice} 和 \texttt{Quot.sound} 三个标准公理。定理、证明及逐项核验记录见形式化附件 \cite{ThisLean}。

对 \(G\) 本身，我们使用 Alpöge 的 Lean 4.33.1 证明与 Sutherland 发布的核验文件 \cite{GRH}。本次独立重放核对了压缩文件及解压文件的 SHA-256，再用 con-leche 的验证模式检查完整证明中的 1\,468\,418 项声明，并与独立命题文件逐项比较。五个目标均由证明给出，其中一个恰为本原狄利克雷特征在开临界带中的 GRH。证明导出只有上述三个标准公理。

谱证明与 GRH 证明分别在各自的版本环境中核验。对导子整除 \(N\) 的特征作特化，即得本文使用的 \(\GRH(N)\)。展开 GRH 命题的名称和数码 \(2\) 的辅助证明，再消去绑定变量名，两份导出中的显式 GRH 命题逐项相同。另逐项比对了狄利克雷特征、本原性、导子、经过较小模数的分解、两个 \(L\) 函数定义及取实部运算。

\begin{thebibliography}{9}
\bibitem{LL} J. Li and J. Liu, \textit{A Conditional Proof of the Goldbach Conjecture},
13 November 2026, 尤见定义 (1.5)--(1.7)、定理 6.1 与第 7 节。
\href{https://mrc.sdu.edu.cn/english/info/1018/1514.htm}{山东大学数学研究中心原文及公告}。
\bibitem{YL} Y. Liu, \textit{Failure of the Li--Liu spectral hypothesis R under Dirichlet GRH},
21 November 2026. \href{https://hexagonmath.org/2611.00788v1}{Hexagon 2611.00788v1}。
\bibitem{LLLean} \textit{Conditional Goldbach in Lean},
\href{https://github.com/subfish-zhou/conditional-goldbach-lean}{形式化定义与证明}。
\bibitem{GRH} L. Alpöge, \href{https://alpo.ge/remark.pdf}{关于黎曼猜想证明的公告}, 20 November 2026.
A. V. Sutherland, \href{https://math.mit.edu/~drew/grh/}{\textit{Verification files for the Lean proof of the Riemann Hypothesis}},
24 November 2026，第五个经核对的命题。
\bibitem{ThisLean} Leiming Fu, \textit{The Li--Liu spectral hypotheses in Lean}, 26 November 2026.
\href{https://github.com/RaymondVoldemortFu/li-liu-spectral-hypotheses-lean}{谱证明与核验记录}。
\end{thebibliography}
\end{document}

